\documentclass[10pt]{article}
\usepackage[T1]{fontenc}
\usepackage{lmodern,amsmath,amssymb,amsthm,microtype}
\usepackage[margin=0.85in]{geometry}
\usepackage[colorlinks=true,urlcolor=blue]{hyperref}
\newtheorem{theorem}{Theorem}
\newtheorem{lemma}{Lemma}
\title{An Exact Denominator Formula for OEIS A334958}
\author{Qichao Wang\\{\small Hebei University of Technology}\\[0.3em]
Daoyu Dong\\{\small University of Electronic Science and Technology of China}}
\date{October 26, 2026}
\begin{document}
\maketitle
\begin{abstract}
We prove the valuation-product conjecture for the gcd of consecutive factorial-scaled alternating harmonic sums in OEIS A334958. An exact reduced-denominator formula proves all of the displayed gcd identities simultaneously and establishes the accompanying divisibility observation. Separate infinitude observations are left open.
\end{abstract}
\section{An exact denominator formula}
For $n\ge1$ put
\[
 H'_n=\sum_{k=1}^n\frac{(-1)^{k+1}}k,\qquad A_n=n!H'_n,\qquad B_n=n!-A_n.
\]
Both $A_n$ and $B_n$ are integers. The sequence in A334958~\cite{oeis} is $a(n)=\gcd(A_n,A_{n+1})$.
\begin{theorem}
Write $H'_n=u_n/v_n$ in lowest terms with $v_n>0$. Then for every $n\ge1$,
\begin{align*}
 a(n)&=\frac{n!}{v_n}
 =n!\prod_{p\le n}p^{\min(0,v_p(H'_n))}\\
 &=\gcd(A_n,n!)=\gcd(B_n,n!)=\gcd(A_n,B_n)
 =\gcd(B_n,B_{n+1}),
\end{align*}
where the product is over primes. In addition,
\[
 \frac{n!}{\operatorname{lcm}(1,\ldots,n)}\ \mid\ a(n).
\]
\end{theorem}
\begin{proof}
The defining sums give
\[
 A_{n+1}=(n+1)A_n+(-1)^{n+2}n!.
\]
Consequently $\gcd(A_n,A_{n+1})=\gcd(A_n,n!)$.
As $A_n=n!u_n/v_n$ is integral and $\gcd(u_n,v_n)=1$, we have $v_n\mid n!$ and
\[
 \gcd(A_n,n!)=\frac{n!}{v_n}\gcd(u_n,v_n)=\frac{n!}{v_n}.
\]
All primes of $v_n$ are at most $n$, and
$\min(0,v_p(H'_n))=-v_p(v_n)$. This proves the valuation-product formula exactly as conjectured in the entry.

The relation $B_n=n!-A_n$ gives the remaining same-index gcd equalities. Also
\[
 B_{n+1}=(n+1)B_n-(-1)^{n+2}n!,
\]
which proves the consecutive-$B$ equality by the same Euclidean step.

Finally the common denominator $L_n=\operatorname{lcm}(1,\ldots,n)$ makes $L_n H'_n$ integral, so $v_n\mid L_n$. Thus
$a(n)=(n!/L_n)(L_n/v_n)$, proving the divisibility observation involving A025527~\cite{lcm}.
\end{proof}
\section{Scope and reproducibility}
The standard-library verifier compares all 29 displayed official terms, independently computes the alternating sum as an exact rational number, and checks the recurrence, denominator, prime-valuation product, every displayed gcd equality, and the divisibility formula for $1\le n\le500$. The included table gives 100 exact terms. Integer sizes grow rapidly and no bit-polynomial complexity claim is made.

The entry also suggests that equality with A025527 and an associated ratio identity occur infinitely often. Neither infinitude assertion is proved here. The completed problem is the explicit all-index valuation-product conjecture, together with the stated divisibility observation; the result is not represented as resolving every question attached to this sequence. No OEIS submission or literature-priority claim is made.
\begin{thebibliography}{2}
\bibitem{oeis} OEIS Foundation Inc., A334958, \url{https://oeis.org/A334958}, accessed October 26, 2026.
\bibitem{lcm} OEIS Foundation Inc., A025527, \url{https://oeis.org/A025527}, accessed October 26, 2026.
\end{thebibliography}

\end{document}
